\chapter{仿真与实验方法}
\section{数学与物理过程分析原理}
\subsection{数字信号处理}
时域分析

对于给定信号s(t)，其时域波形的基本特征常用信号能量、波形中心、均方持续及时宽等特征量来描述。

频域分析

信号的频域描述就是利用傅里叶变换，将信号 s(t) 用不同频率的正弦波展开。即信号由简单的正弦波相加（线性叠加）组成，每一个正弦波都有一个频率和用系数表示的有关的量。

时频分析

然傅里叶变换（FT）有很大的优越性，但它却存在许多不足。首先，在以外的空间，虽然傅里叶变换式在一定条件下仍然成立, 但大小不能刻画出s(t)所在的空间。其次，用傅里叶变换只能获得信号s(t)在时间范围内的频谱，而难以了解信号在某时间范围内的性质，即它对频率的分辨率是无穷的，而对时间的分辨则为零，不能同时对时间和频率具有较好的分辨率，这是因为傅里叶变换是对变量t求积分，去掉了非平稳信号中的时变信号，它只适宜于确定性的平稳信号，而对时变非平稳信号则难以充分刻划。同时傅里叶变换对时域的分辨率在任一区间上也是不变的，因而不足以在任意小的范围内描述或确定信号s(t)。

然而，对于该研究下的时变频率信号，我们关心的恰恰是信号在局部范围内的特性。如什么位置出现什么样的异常信号、信号突变的位置、信号在突变时刻所对应的频率成分，显然傅里叶变换的积分作用平滑了非平稳过程的突变成分，作为积分核的 的幅值在任何情况下均为 1，丢失了需要研究的信号。

时频分析（time-frequency analysis, TFA）是非平稳信号分析最基本的内容， 其基本任务就是建立一个以时间t和频率为变量的二维联合分布函数，简称时频分布。利用该分布函数可以求某一确定的频率和时间范围的能量分布，计算在某一特定时间的频率的密度，计算该分布的整体和局部各阶矩 （moment），如平均条件频率及其局部扩展等，利用这种时频分布来讨论非平稳 信号时频变特性，如瞬时频率、瞬时带宽、群延迟等。

\subsection{机器学习}
\subsection{仿真建模}

\section{仿真建模}
% TODO: 草稿中“仿真建模”标题下仍缺少正文。请补充 COMSOL/等比模型/缩比模型的建模流程、参数和边界条件。
\section{实验系统与数据采集}
% TODO: 请补充 MEMS 传感阵列、水槽/水下实验平台、采样频率、通道定义、目标运动方式和原始数据格式。
